% =============================================================================
%  Principles of Robot Motion via Rust — The FCP Way
%  Print edition, in the Shannon Robotics design language.
%
%  Build:  make        (or: pdflatex main.tex, three times)
% =============================================================================
\documentclass[11pt, twoside, openright]{book}

\usepackage[
  paperwidth=7in, paperheight=10in,
  inner=0.95in, outer=0.75in, top=0.85in, bottom=0.95in,
  headsep=16pt, footskip=28pt
]{geometry}

\usepackage{shannon}

\hypersetup{
  pdftitle={Principles of Robot Motion via Rust --- The FCP Way},
  pdfauthor={Shannon Robotics},
  pdfsubject={Robot motion planning, configuration space, sampling-based planning, nonholonomic systems, Rust},
  colorlinks=true,
  linkcolor=shAccent,
  urlcolor=shSkyDark,
  citecolor=shAccent
}

\begin{document}

% =============================================================================
%  Front matter
% =============================================================================
\frontmatter
\pagestyle{empty}

\begin{titlepage}
\thispagestyle{empty}
\noindent
{\color{shTeal}\rule{\textwidth}{3pt}}
\vspace{2.2in}

\noindent
{\displayfont\fontsize{10}{12}\selectfont\color{shTealDark}%
 \MakeUppercase{Shannon Robotics}}

\vspace{10pt}
\noindent
{\displayfont\bfseries\fontsize{30}{35}\selectfont\color{shInk}%
 Principles of\\Robot Motion\\via Rust}

\vspace{12pt}
\noindent
{\displayfont\fontsize{16}{20}\selectfont\color{shInkMuted}%
 The FCP Way}

\vspace{20pt}
\noindent
{\color{shSlate300}\rule{2.4in}{0.6pt}}

\vspace{16pt}
\noindent
\begin{minipage}{4.2in}
{\fontsize{11}{15}\selectfont\color{shInkSecondary}%
 Complete mathematical foundations, interactive intuition, and working
 implementations in Rust --- for robots that know where they are and must
 decide how to move.}
\end{minipage}

\vfill
\noindent
{\displayfont\fontsize{9}{12}\selectfont\color{shInkMuted}%
 Foundation \ \textperiodcentered\ \ Conceptual \ \textperiodcentered\ \ Practical}

\vspace{6pt}
\noindent
{\color{shTeal}\rule{\textwidth}{3pt}}
\end{titlepage}

% ----- Colophon --------------------------------------------------------------
\cleardoublepage
\thispagestyle{empty}
\vspace*{\fill}
\noindent
{\fontsize{9}{13}\selectfont\color{shInkSecondary}
\textbf{Principles of Robot Motion via Rust --- The FCP Way}\par
\vspace{5pt}
Print edition, generated from the interactive web edition.\par
\vspace{9pt}
This book takes Choset, Lynch, Hutchinson, Kantor, Burgard, Kavraki and Thrun's
\emph{Principles of Robot Motion} as its spine and brings it forward:
asymptotically optimal sampling-based planners, trajectory optimization and
model predictive control, multi-robot and manipulation planning, learning inside
the planning loop, and an integrated planning stack. It is the sister volume of
\emph{Probabilistic Robotics via Rust}, with which it shares its robot, its
simulator and its method; estimation is treated here only as far as a planner
needs it, and cross-referenced there for the rest.\par
\vspace{9pt}
Every chapter of the web edition carries interactive simulations that run the
same algorithms the text derives. In this print edition each appears as a
described figure with its widget identifier, so the two editions can be read
side by side.\par
\vspace{9pt}
{\color{shInkMuted}Set in Space Grotesk and a Palatino-class serif.
Composed with \LaTeX.}\par}
\vspace*{\fill}

% ----- How to read this book -------------------------------------------------
\cleardoublepage
\pagestyle{fancy}
\chapter*{How to read this book}
\addcontentsline{toc}{chapter}{How to read this book}

Every chapter makes three passes over its subject, and they are meant to be read
in this order.

\textbf{Conceptual} comes first, because intuition should precede formalism. In
the web edition you meet an interactive figure before you meet an equation; here
that figure is described, and the description tells you what it would show you.

\textbf{Foundation} is the mathematics, in full. Definitions are precise,
assumptions are stated where they are made rather than buried, and derivations
are carried to the end. Long algebra sits in marked derivation blocks, so you
can take the result on faith the first time through and come back for the proof
later.

\textbf{Practical} is the Rust. Not pseudocode dressed up as code: real types,
real crates, code you could lift into a robot.

\section*{The running lab}

Three robots and three worlds recur throughout, so that every new method can be
compared against the last on identical ground. \shrobotname{Rusty} is a
differential-drive rover with a range sensor --- the same robot as in the sister
volume --- and lives in the \textbf{Apartment}, a 2-D floorplan with rooms,
doorways and a long corridor; for most of Parts~I--IV it is treated as a disc.
\shrobotname{Reach} is a planar two-link arm (with a three-link variant) working
at the \textbf{Workbench}, a table with polygonal obstacles; its configuration
space is a torus, which is why it appears whenever topology matters, and its
dynamics carry Part~V. \shrobotname{Hitch} is a car-like robot that can pull a
trailer and lives in the \textbf{Lot}; it cannot move sideways, which is the
whole subject of Chapters~20 and~21. All three are built in Chapter~2 and act
together in Chapter~23.

\section*{Colour carries meaning}

One convention runs through the prose, the equations, the figures and the code
comments alike. It is worth internalising before you start:

\shplanningkey

\noindent
Green is always where you start and red is always where you want to be. Slate is
a thing in the world; amber is the set of configurations that would make the
robot touch it. Blue is what a planner has explored --- a roadmap, a tree, a
wavefront. Purple is the path it chose. Orange is the robot itself and the
trajectory it actually executed. Where a belief appears (Chapters~15--16 and the
capstone) the sister volume's five estimation colours are used unchanged:

\shcolorkey

\noindent
The teal used for chapter furniture is deliberately outside both sets, so it can
never be mistaken for data.

% ----- Contents --------------------------------------------------------------
\cleardoublepage
\setcounter{tocdepth}{0}
\renewcommand{\contentsname}{Contents}
\tableofcontents

% =============================================================================
%  Main matter
% =============================================================================
\mainmatter

\part{Foundations: Robots, Worlds, and Configuration Space}
\input{chapters/ch01}
\input{chapters/ch02}
\input{chapters/ch03}
\input{chapters/ch04}
\input{chapters/ch05}

\part{Classical Planners: Search, Potentials, Roadmaps, Cells}
\input{chapters/ch06}
\input{chapters/ch07}
\input{chapters/ch08}
\input{chapters/ch09}
\input{chapters/ch10}

\part{Sampling-Based Planning}
\input{chapters/ch11}
\input{chapters/ch12}
\input{chapters/ch13}
\input{chapters/ch14}

\part{Planning with Uncertainty}
\input{chapters/ch15}
\input{chapters/ch16}

\part{Dynamics, Trajectories, and Constraints}
\input{chapters/ch17}
\input{chapters/ch18}
\input{chapters/ch19}
\input{chapters/ch20}
\input{chapters/ch21}

\part{Frontiers and Integration}
\input{chapters/ch22}
\input{chapters/ch23}

\end{document}
